\par\noindent\textbf{(1)} 体坐标对矢量求导满足 (dL/dt)\_space=(dL/dt)\_body+ω×L。无外力矩、L\_i=I\_iω\_i，得到 I1ω̇1=(I2−I3)ω2ω3、I2ω̇2=(I3−I1)ω3ω1、I3ω̇3=(I1−I2)ω1ω2。乘各ω\_i求和，交叉项抵消，故 E=ΣI\_iω\_i²/2 不变；同理 L²=ΣI\_i²ω\_i² 不变。
\par\noindent\textbf{(2)} 绕1轴令ω1=Ω、ω2=u、ω3=v，小扰动方程 u̇=(I3−I1)Ωv/I2、v̇=(I1−I2)Ωu/I3。因此 ü=−[(I3−I1)(I2−I1)/(I2I3)]Ω²u，角频率 ν1=|Ω|√[(I3−I1)(I2−I1)/(I2I3)]，两正因子保证中性小振动稳定。
\par\noindent\textbf{(3)} 绕2轴设ω2=Ω，ω1=u、ω3=v，得 u̇=(I2−I3)Ωv/I1、v̇=(I1−I2)Ωu/I3；两系数同为负，故 ü=λ²u，λ=|Ω|√[(I3−I2)(I2−I1)/(I1I3)]>0。一分支按e\textasciicircum{}\{λt\}增长，另一按e\textasciicircum{}\{−λt\}衰减，这是网球拍翻转的线性机制。
\par\noindent\textbf{(4)} 绕3轴的类似计算给 ν3=|Ω|√[(I3−I2)(I3−I1)/(I1I2)]，因此最大、最小主惯量轴都线性稳定。若I1=I2或I2=I3，某频率或增长率趋0，需重新考虑轴对称刚体的连续退化模，不能用严格不等式下的稳定类型直接外推。
